% !TEX root = ../main.tex
\chapter{Newtonian dynamics (Newtonsche Dynamik)}\label{ch:c10}
\begin{paracol}{2}

\begin{sources}
\begin{tabularx}{\linewidth}{@{}l l L@{}}
\toprule
\textbf{Video} & \textbf{Length} & \textbf{Content}\\
\midrule
\seg{10}{1}{Lecture 10: Newtonian dynamics} & 1:44:53 & Why the force concept is strange and the field idea (Faraday); spatial vector fields and the interaction fields (Wechselwirkungsfelder), the force on a particle;
why the first axiom speaks of vanishing charges; which interaction fields exist (fundamental ones, effective ones) and why electromagnetism does not fit the Newtonian spacetime;
outlook to the last lecture; the precise second Newtonian axiom, its chart form and the textbook special case; the problem sheet; work, conservative interaction fields and potential energy.\\
\bottomrule
\end{tabularx}

\smallskip
\textbf{Transcript} \file{materials/transcripts/C10-V01.txt} (German YouTube auto-captions; good quality, a few garbled formula dictations are reconstructed from context and marked [?] where unsure).
\end{sources}

\begin{reading}[Where this lecture sits in the other notes]
\emph{Gravity and Light}: lecture 9, ``Newtonian spacetime is curved'' (there the dynamics, the force written as a field and the gravitational force as a component of the connection are put together: the gravitational field is
traded for the curvature of the Newtonian connection, which is announced here and in the exercise on the curved Newtonian spacetime). Continues \cref{ch:c09} (spatial spacetime metric and absolute space, which are needed for work and potential energy);
the next lecture (\cref{ch:c11}) specialises to the gravitational interaction field.
\end{reading}

% ==================================================================
\section{Forces, fields and the plan of the lecture}\label{sec:nd-plan}

\begin{aside}[Organisation: the date of the exam]
Before the lecture the lecturer asks the audience about the date of the (mock) exam: the date he had proposed, Saturday the 7th, is unlucky because of the long weekend; the only alternative is two Saturdays later, the 21st. A show of hands (who
can on the 21st, who cannot, who can on the 7th, who cannot) gives a split audience; he will think a day or two, perhaps make a Doodle poll; ``to me it is all the same, the 7th or the 21st, all other dates I am booked'' \ts{10}{1}{0:09}.
(Not part of the content; kept for completeness.)
\end{aside}

We start with the tenth lecture, and finally come to the Newtonian dynamics, which is, roughly speaking, about the \emph{second Newtonian axiom} and thereby about \emph{forces} \ts{10}{1}{1:12}.
The concept of force is familiar from school, but as it occurs in Newton it is something incredibly strange \ts{10}{1}{1:32}. There is talk of the force that one particle exerts on another one (the third axiom has it explicitly:
equal and opposite forces and so on). But what does this mean? If one particle exerts a force on another, how does the force get from one particle to the other? How does this interaction take place \ts{10}{1}{1:53}?

\paragraph{Instantaneous action contradicts relativity.} You all watch the news and know there is a general theory of relativity, in which the propagation of signals (the speed with which one can transmit information) is bounded by the speed of light.
Suppose I hold one particle and you hold a similar one, and my particle exerts a force on yours, as in the Newtonian picture. Now I wiggle my particle; then the force vector acting on your particle would, by this philosophy, align
itself \emph{instantaneously} with the new position (think of gravitation or the electric force). If that were the case, and nothing else is said in Newtonian mechanics, I could send you, no matter how far away you are, Morse signals
instantly (``three short, three long, three short'') \ts{10}{1}{2:15}. This cannot be so in relativity, and it is not \ts{10}{1}{3:17}.
One could say that this does not interest us, since we do non-relativistic mechanics. But electrodynamics, which you will meet next semester, was written down with the kinematics that we have developed so far for Newtonian mechanics in the back of the head:
Maxwell wrote it essentially from the Faraday and the Amp\`ere laws, with a small additional term fiddled in so that the conservation of charge is guaranteed; it \emph{looks} like a beautiful non-relativistic theory, but it is already relativistic, one merely does not
notice at first sight; it took Einstein to look longer at it \ts{10}{1}{3:39}.

\paragraph{The field concept (Faraday).} Faraday was one of the persons to whom a whole lecture of next semester is devoted; he could not do mathematics, and therefore he invented the \emph{field} concept: one has drawn it in school, from a point charge go field lines, and there
is this force. Astonishingly this idea, that the force does not simply act but acts \emph{through a field}, is the one we have today: a particle \emph{creates} a field and the other particle \emph{reacts} to this field in a certain way; we may think that the field produces
at the point where the other particle sits the force \ts{10}{1}{4:22}. This notion of field, created by one particle and reacted to by another, appears nowhere in Newton; it is an idea developed much later by Faraday, and it is the idea we have modern; therefore the
chapter title ``Newtonian dynamics'' is conceptually, sensibly, somewhat \emph{modernised} \ts{10}{1}{5:25}. A little classification, so that the concepts do not stand unconnected side by side \ts{10}{1}{5:47}.

\begin{coursenote}[Faraday, Maxwell, Einstein: what was said and what is not checked]
Everything in this section about Faraday, the Faraday and Amp\`ere laws, Maxwell's extra term and Einstein ``looking longer'' is as said in class (not checked); the lecture only gives these facts in passing, as motivation for the field concept and for the failure of the Newtonian
framework for electrodynamics discussed below (\cref{sec:nd-which}).
\end{coursenote}

% ==================================================================
\section{Interaction fields and the force on a particle}\label{sec:nd-fields}

First, the interaction fields (Wechselwirkungsfelder) \ts{10}{1}{5:47}. Two points are recalled from the last lectures \ts{10}{1}{6:12}. The first is a technical one; the second becomes relevant at once.

\begin{recall}[Spatial vector fields and particles (\cref{ch:c09,ch:c08})]
A \emph{spacetime} vector field is called spatial if its component functions in a chart of a stratified atlas of the Newtonian spacetime have the form in \eqref{eq:nd-spatial} below (the action of the projection $\delta^\alpha{}_a$ of \cref{def:as-G});
a \emph{particle} is a tuple $(\gamma,m,c_1,\dots,c_L)$ of a future-pointing world line, an inertial mass $m>0$ and $L$ charges (\cref{def:nk-particle}).
\end{recall}
The first point: a \textbf{spatial spacetime vector field} (the emphasis again on which space the field lives, because we now have the absolute space as well: a ``space vector field'' would be a field on the space, but we consider first a vector field on the whole spacetime)
\ts{10}{1}{6:12}. Notice that it is still a spacetime vector field, but it is called \emph{spatial} if its component functions at a point $p$ are given by
\begin{equation}\label{eq:nd-spatial}
  X^a(p)\;=\;X^\alpha(p)\,\delta^\alpha{}_a\qquad[\text{i.e.\ } X^0=0,\ X^a=X^\alpha\text{ for }a=\alpha=1,2,3],
\end{equation}
which is what we talked about the other day \ts{10}{1}{6:54}. This is \emph{not} to hold in an arbitrary atlas but only in a chart of a \textbf{stratified atlas} of a Newtonian spacetime (``I do not write this each time, you know the context'') \ts{10}{1}{7:37}.
(The index placement on the board, ``$X^\alpha$ and at the back a $\delta^\alpha{}_a$'', is garbled in the transcript; the form \eqref{eq:nd-spatial} is the one of \cref{ch:c09}, so that it just says: zero time component; [?].)

The second point, which will become very relevant: the definition of a particle, a tuple consisting of a future-pointing world line, a positive real number called the inertial or ``heavy'' mass and a number $L$ of charges \ts{10}{1}{8:05}.
Today we are particularly concerned with the charges. About them nothing much has been said so far, except that in the first axiom we looked only at particles all of whose charges vanish. This is already the modern view: in Newton there is no talk of charges,
whereas of the mass of a particle there is. We add them \ts{10}{1}{8:48}. The decisive point: to these $L$ charges there correspond $L$ \textbf{interaction fields} \ts{10}{1}{9:08}.

What value does $L$ take? It depends on the universe one considers \ts{10}{1}{9:51}. You have probably heard that there are four fundamental forces: gravitation, electromagnetism, the weak force and the strong force. Originally there were gravitation and electric and magnetic force,
and Maxwell united the electric and the magnetic force into the electromagnetic one; then (in the 70s, ``Weinberg, Glashow, Salam'') electromagnetism was united with the weak force into the \emph{electroweak} force; so even $E$ and $B$ together are not yet a complete theory,
and much work in physics of the last decades went into attempts to unite the strong force too. This is the idea of the unification of forces: some people pursue the aim to make $L$ smaller and smaller, to have fewer and fewer fundamental interaction fields that
produce the forces \ts{10}{1}{10:12}. But we do not want to commit ourselves: $L$ depends on the universe considered. One may imagine our universe with a \emph{fifth} force (people have looked for one; we do not know; perhaps someone finds it, or one finds that the forces
are not independent). ``$L$ depends on the universe'' in the sense that the theory can deal with any number $L$ \ts{10}{1}{11:14}.

\begin{definition}[New concepts: interaction field (Wechselwirkungsfeld)]\label{def:nd-field}
An \textbf{interaction field} is a spatial spacetime vector field $W$ (numbered $W_1,\dots,W_L$, one for each charge $c_i$) \ts{10}{1}{11:55}.
\end{definition}
This is not yet particularly enlightening. The remark: the idea behind the \emph{fiction} of interaction fields is the following \ts{10}{1}{12:44}. The fields could just be \emph{there}: a Newtonian universe could consist of very many point particles and of interaction fields thrown in somewhere, simply
present. One can do that, and one often computes like that: one knows these exercises, ``consider a charged particle in an electric field''. Where does the electric field come from? It is simply there; one ignores where it comes from. If one thinks the thing completely fundamentally,
we assume today that the fields, to which the particles react, are \emph{themselves produced by the particles}: a particle does not only react to a field if it has a charge, it also produces such a field \ts{10}{1}{13:16}. The question how a distribution of particles produces one of these fields
is answered by \textbf{field equations}: all particles of the universe together produce the $L$ interaction fields by means of suitable field equations (``such as \dots{} later, I cannot say it yet, it will come'') \ts{10}{1}{14:17}. The picture is clear: a particle produces a field and
the others react to it, but the reacting particle also contributes to the field. So there are $L$ interaction fields in our universe, \emph{not} a new vector field for each particle: all particles have charges $c_1,\dots,c_L$ and there are really only $L$ interaction fields \ts{10}{1}{15:26}.
Imagine $c_1$ a gravitational charge (``although I always preach that it is not a force, let us assume it for a moment'') and $c_2$ an electric charge; then $W_1$, the interaction field belonging to the charge $c_1$, would be something like the electric field and $W_2$ something like the
gravitational field (the lecture mixes the labels; the pairing of the numbers is irrelevant) \ts{10}{1}{16:07}.

\begin{definition}[New concepts: the force on a particle (Kraft)]\label{def:nd-force}
The \textbf{force} acting on a particle $(\gamma,m,c_1,\dots,c_L)$ is the spatial Newtonian spacetime vector field
\begin{equation}\label{eq:nd-force}
  F\;=\;c_1W_1+c_2W_2+\dots+c_LW_L ,
\end{equation}
where the $c_i$ are the charges of this specific particle \ts{10}{1}{16:29}.
\end{definition}
So the $W_1,\dots,W_L$ are there because they were somehow produced, whereas $c_1,\dots,c_L$ are the charges of \emph{this specific particle}; and this we call the force on the particle \ts{10}{1}{17:15}. Suppose a particle with certain charges sits at the tip of my finger, and everywhere there are
the interaction fields (vector fields, so everywhere). One can of course form the sum \eqref{eq:nd-force} also at a point where the particle is not; one could call it the force there, but this brings nothing: it is the force for exactly this particle, so it makes sense only where the particle is.
In principle $F$ is defined everywhere on the Newtonian spacetime, but it matters only at the points $\gamma(\lambda)$, $\lambda\in\R$, those past which the curve comes \ts{10}{1}{17:56}. Another particle, at another position or the same one, feels a \emph{different} force, since its charges can
differ; but the interaction fields are independent of the particle on which they act \ts{10}{1}{19:00}. ``You know all these concepts, but the Newtonian force always composes itself like this, in the modern view'' \ts{10}{1}{19:20}.

\paragraph{Why the first axiom speaks of vanishing charges.} Now it is finally clear why in the first axiom we did not say ``particles on which no force acts \dots'' but ``particles whose charges all vanish'' \ts{10}{1}{19:41}. Vanishing charges certainly lead to no force; but I can \emph{test beforehand} whether a particle has no charge:
I take the particle into the laboratory, subject it to all possible tests (swing magnets, move large masses back and forth), and if I find no way of attributing any charge to it, because it does not react to anything I do, then I say it has zero charge with respect to the interaction field.
Then I can take this particle, recognised as chargeless, and do my throwing experiment far out in space and say: ah, these are straight lines. The particle charge as a property of the particle is needed to procure proper matter of which I can later say that no force acts on it; otherwise I would not know
that no force acts: the cloud of charges makes this more transparent \ts{10}{1}{20:01}. If there are no interaction fields at all, no force would act either \ts{10}{1}{21:24}.

(No question from the audience at this point: ``is this all clear so far, or are there questions? Good'' \ts{10}{1}{21:24}.)

% ==================================================================
\section{Which interaction fields exist? Fundamental and effective ones}\label{sec:nd-which}

Now, since this is physics and not mathematics, we must show our colours and say which interaction fields there \emph{are} in nature \ts{10}{1}{21:44}. Here too it pays to be terminologically exact. The interaction fields fall, by decree, not as a natural observation but owing to a love of
order, into two classes \ts{10}{1}{22:09}.

\subsection{Fundamental interactions}\label{ssec:nd-fundamental}

The first class: \textbf{fundamental interactions} (fundamentale Wechselwirkungen). ``Fundamental'' is a word easily misused: in politics one says ``fundamentalist'' and polemicises, in physics one says ``fundamental'' and means something great; actually it does not mean much,
except that these are interactions that we have so far not been able to break up into still more elementary interactions \ts{10}{1}{22:30}. The list \ts{10}{1}{23:21}:
\begin{enumerate}
  \item gravitation (the corresponding gravitational interaction field);
  \item the electric field;
  \item the magnetic field;
  \item the weak (nuclear) force, ``weak field'' (there is no particularly illustrative name since one has no intuition for it on the macroscopic level);
  \item the strong (nuclear) force, ``strong field''.
\end{enumerate}
Since Maxwell the electric and the magnetic field are understood as the \emph{electromagnetic field}, and today it is even united with the weak one in the \emph{electroweak interaction} (the magnetic then falls under the table, but it swings along); the strong one stands for the moment
a bit alone. This is the unification process which has worked twice, from which the physicist likes to conclude that it must continue forever, ``for which there is of course no logical reason: it is pure principle of hope; about 30 years of research went into this without effect, one has to say''
\ts{10}{1}{24:26}.

The weak and the strong field (one feels it: beta decay and such) have no proper classical description at all; for this one needs quantum mechanics, which one learns in the third theory lecture, or next semester for teaching-degree or materials-physics students in connection with the electromagnetic field. These two certainly
do not fit into classical mechanics: there is no use in trying to write them down as interaction fields in classical mechanics, ``we have no chance'' (the lecturer groups them as \emph{not classically describable}, the others as \emph{classically describable}) \ts{10}{1}{25:49}.

\paragraph{Gravitation.} Gravitation is super: it is classically describable, even non-relativistically, all fine, except for the point (it is hanging out of your throat already) that we do not want to understand gravitation as a force and an interaction but as the \emph{curvature of spacetime}. We
really come to this in the \emph{next lecture} \ts{10}{1}{26:31}. It can be described as an interaction in the Newtonian spacetime, but ``we do not want to''; we shall nevertheless do it, simply since one can compute so cleverly with it; we do not want it as a fundamental idea
\ts{10}{1}{27:12}.

\paragraph{The electric and the magnetic field do not fit the Newtonian spacetime.} What remains are the electric and the magnetic field, and here there is a bad piece of news that is easy to see \ts{10}{1}{27:57}. If the electric field were an interaction field it would have to be a spatial spacetime
vector field; one thinks of $\vec E$, which has the little arrow on top, that is, a three-dimensional vector; but you know better by now, it has components $E^\alpha$ and is to transform as a vector. We have a little problem, as follows
(thought experiment, as said in class) \ts{10}{1}{28:19}.
\begin{enumerate}
  \item A conductor (a wire) which is charged, but no current flows; the electrons sit evenly distributed and produce the cylindrically radial electric field: \emph{an $E$ field and no magnetic field} \ts{10}{1}{28:40}.
  \item Now let a uniform current flow through the conductor: the electrons move; the magnetic field is the circular one around the wire: \emph{a $B$ field}. In the first case one has an $E$ field, in the second a $B$ field, apparently with nothing to do with each other \ts{10}{1}{29:01}.
  \item Return to the first thought experiment, with only the $E$ field, the charges at rest. Now go to a \emph{different inertial system}: ``I stroll past the conductor''; this inertial system is as good as any other, but in it the charges \emph{move}. And if the charges move, I must see a magnetic field,
        correct? And so it is: one sees this $B$ field \ts{10}{1}{29:21}.
\end{enumerate}
So one has $E$ and $B$ field, apparently independent; one has an $E$ field and no $B$ field, goes to another inertial system, and suddenly a $B$ field grows. ``That is crazy; yes, right, it is crazy, and this is electrodynamics, and it has nothing to do with non-relativistic mechanics''. $E$ and $B$ are not two
spacetime vector fields but \emph{both components of one spacetime $(0,2)$ tensor field}; only this tensor field transforms under the corresponding coordinate transformations, and if one looks closely they are not the transformations that occur in our Newtonian framework but those of relativity
\ts{10}{1}{30:03}. In short: electric and magnetic are two sides of one coin, they cannot be separated \ts{10}{1}{30:44}. Staying always in one inertial system one may permit oneself the fantasy ``there is an $E$ field and a $B$ field, and now we compute with Newtonian mechanics how the particles move''; in principle this is
also right but ``conceptually it is complete nonsense'' \ts{10}{1}{31:05}. The only one who recognised this, after (the lecturer corrects himself) 50 years of Maxwell theory, was Einstein: he looked and said, that is not right \ts{10}{1}{31:25}. We shall do the same at the end (``in the last but one lecture or so'') and do it better;
the problem is for the moment not describable in the Newtonian spacetime: classically describable, but not in the Newtonian spacetime \ts{10}{1}{31:48}.

Consequently the sobering conclusion: of all fundamental interactions experimentally established until today, if one takes it very precisely and does \emph{not} want to describe gravitation as an interaction (in order not to make the first axiom unemployed), \textbf{no force is left} \ts{10}{1}{31:48}. ``So it is. So it is.''

\begin{coursenote}[The $E$/$B$ argument is a motivation, not a proof]
The thought experiment is told qualitatively (no formulas, no transformation law is derived); the claim that $E,B$ form the components of a spacetime $(0,2)$ tensor field is stated as a fact of electrodynamics, to be seen next semester, and the lecture promises to redo it in the relativistic
setting at the end of the semester (\cref{ch:c22,ch:c23}). Nothing is proved here, and none of it is used below.
\end{coursenote}

\subsection{Effective interactions}\label{ssec:nd-effective}

But in experimental physics you have looked for hours at force fields and computed with them, potentials, the harmonic oscillator and all kinds of knick-knack. What is that? It is a concept we introduce now: these are not fundamental but \emph{effective} interactions \ts{10}{1}{32:30}.
This is the \textbf{second type of interaction field: effective interactions and interaction fields} \ts{10}{1}{32:54}.

What is it? The physicist is a practical person by nature, and he says: my physical thinking I can apply to anything that can be thought of in nature, I analyse it to death \ts{10}{1}{33:17}. Now one may observe interactions of some objects, for example of human beings. A human is very complicated: it contains the strong and the weak and the
magnetic force and gravitation and all kinds of things. (Our circulation, that the heart has to pump and how it pumps, all has to do with our being in the gravitational field; if you fly into space you must take special precautions for your circulation. To understand a human one must understand at least these five forces; if one is a
reductionist one then says that chemistry and our thinking and, out of thinking, our actions, all follow from these \ts{10}{1}{33:38}.) But if that is so, the idea of \emph{effective interaction fields} is this \ts{10}{1}{34:19}.
\begin{example}[An effective interaction between two human beings]\label{ex:nd-man}
(``Normally I do not like to make examples, but for effective interaction fields there are only examples'' \ts{10}{1}{34:39}.) Consider the interaction of one particle (a man) with another (a ``sausage'': in the lecture a more coloured pair of persons; the point is the same).
These are two very complicated objects, with chemistry and all that rubbish, but we do not want to understand the interaction fundamentally: we want to make a \emph{model} \ts{10}{1}{35:00}. The observation: if someone were to walk normally (the free path, the green straight lines, the autoparallels, drawn in space)
and there is the other one, then something happens, the path bends to the other; there is something complex (chemical, optical, photons are exchanged first), and there is in any case an \emph{overall effect}, and this overall effect one can \emph{model as an interaction} \ts{10}{1}{36:03}.
It depends of course on the charges: today it is easy to say that all particles carry charges, and one has different kinds of charges (various, the lecturer jokes, even ones ``with a beard''); one sees how \emph{$L$ grows}. One has the fantasy to see in which cases which charge is how (plus, zero, twice plus) \ts{10}{1}{36:46}.
It is an \emph{effective modelling}: we need not understand where the attraction between the man and the sausage comes from; it can simply be modelled effectively. ``But you understand it is nothing fundamental, it comes just out of the hand'' \ts{10}{1}{37:50}.
\end{example}
\begin{example}[The Hooke spring]\label{ex:nd-hooke}
Another example: a wall, a little cart, a mass fixed on the cart, a spring. One learns the Hooke force law and so on. But honestly: when I pull away the cart and stretch the spring, the spring consists of atoms and material; the strong field lets them hold together at all, together with the
electromagnetic one; a spring is an incredibly complicated thing that one understands fundamentally only quantum-mechanically, and even then \ts{10}{1}{38:10}. The restoring force exerted by the spring is only something \emph{effective}: not understood fundamentally. A spring is called a Hooke spring if in the regime in which
one looks at it it has a linear restoring law; and this is nothing fundamental, just ``thrown down'' as a force law (the lecturer's word is rougher). One may then consider non-Hooke springs with a cubic term and so on, ``all only effective'': one looks at what happens and then invents one's force so that, when the second
Newtonian axiom (which comes in a moment) is used, the motion comes out that one sees anyway \ts{10}{1}{39:11}.
\end{example}

\paragraph{Summary: ad hoc interaction fields.} Often one does not understand (or does not want to deal with) \emph{how exactly a deviation of the world line of a particle from the straight lines} is to be traced back to the fundamental interactions; then one models the interaction fields \textbf{ad hoc}. The lecturer
admits that the formulation is a bit ``larifari'' but the point is only to hold the idea \ts{10}{1}{40:15}. One invents charges and interactions; to return to the human example, one can then for instance give the human drugs so that a charge is switched off or particularly switched on: one can manipulate with this, and such things
work \ts{10}{1}{42:14}. Likewise for the other things \ts{10}{1}{42:35}.

\paragraph{Outlook.} Since we have lost almost all fundamental interactions, what remains in mechanics is a bit sad, but it is so; it does not change that the fundamental ones would be the aim, which \emph{Newtonian} mechanics cannot achieve according to the present state of knowledge.
You might think, oh dear, this is really weak, but the point is that we are incredibly close \ts{10}{1}{42:55}. An outlook (``in green, like the hope''), to the \emph{last lecture of this semester} and to the further theory lectures \ts{10}{1}{43:21}.
So far we had $(M,\mathcal O,\mathcal A,\nabla,t)$ and a spatial spacetime metric $G$ (together a Newtonian spacetime with absolute space), and compatibility conditions such as $\nabla G=0$ \ts{10}{1}{43:41}. Essentially it is about \emph{striking out the absolute time}: one defines a new kind of
spacetime in which there is no absolute time any more, therefore no simultaneity (how would it be defined?); then $G$ is not a \emph{spatial} spacetime metric any more either, since the spatial in the spatial spacetime metric was defined by means of the absolute time $t$; one takes a spacetime metric without the
restriction ``spatial'', with essentially the same conditions on the remaining objects, and that is an \textbf{Einstein spacetime} (Einsteinsche Raumzeit) \ts{10}{1}{44:02}. This is general relativity, and once we have written it down, \emph{special} relativity follows
``really only as a small application''. So with the absolute time Newton shot himself in the knee, and Maxwell theory, even if written in this framework, blows up the framework \ts{10}{1}{44:45}. ``Now you perhaps also see why I practise this fanaticism with all these structures: one understands what happens in
the transition from non-relativistic to relativistic physics''. And then, miracle of miracles, \emph{all five forces fit in}; we are now mathematically and philosophically prepared, we only have to get rid of the absolute time, which happens at the very end of the semester; ``so it is half as bad as one would think
after this catastrophic intermediate result'' \ts{10}{1}{45:28}.

% ==================================================================
\section{The second Newtonian axiom}\label{sec:nd-second}

We come to the application of these notions (interaction fields and force): the precise formulation of the \textbf{second Newtonian law} \ts{10}{1}{46:11}.

\begin{theorem}[Second law of motion (zweites Bewegungsgesetz); second Newtonian axiom]\label{ax:nd-second}
The world line $\gamma$ of a particle $(\gamma,m,c_1,\dots,c_L)$ in a universe with interaction fields $W_1,\dots,W_L$ satisfies the differential equation
\begin{equation}\label{eq:nd-second}
  m\,\nabla_{\!X_\gamma}X_\gamma\bigl(\lambda\bigr)\;=\;c_1W_1+\dots+c_LW_L ,
\end{equation}
where $X_\gamma$ is the tangent vector field of $\gamma$, so that the left-hand side is $m$ times the \textbf{acceleration} of the particle (the covariant derivative of the tangent along its own direction, identified in an earlier lecture, \cref{ch:c08}) at the parameter value $\lambda$, and every $W_i$ on the
right-hand side is evaluated at the point $\gamma(\lambda)$ where the particle is. Without loss of generality the parameter is the absolute time of the event: $\lambda=t$ \ts{10}{1}{46:39}--\ts{10}{1}{50:25}.
\end{theorem}
(It is the second axiom; the box type of the notes marks a statement, it is postulated and not proved.) In words \ts{10}{1}{47:13}: the world line of a particle, with $\gamma$ the world line, $m$ the inertial mass and $c_1,\dots,c_L$ the charges that the particles have in the universe we assume, satisfies a
differential equation (as in the first axiom): the object $\nabla_{X_\gamma}X_\gamma$ does not frighten us any more, it is just the \emph{acceleration} of the particle, identified in one of the first lectures; mass times acceleration equals a force, which however we write in modern spirit as the
\emph{weighted sum of the interaction fields with the charges of the particle}, instead of the traditional ``$F$'' \ts{10}{1}{48:19}.

\paragraph{The parametrisation.} If one admits arbitrary parametrisations, an affine reparametrisation $\lambda=at+b$ would make the factors $a$ appear (and the right-hand side is no longer $0$, so the factors cannot be divided out). We therefore make, \emph{without loss of generality}, the choice that the curve parameter
equals the absolute time of the event $\gamma(\lambda)$: ``this is just a choice out of convenience, you can do it without destroying anything'' (what it means: $\lambda$ is a number, $t$ is a function; ``$\lambda=t$'' is formal nonsense but to be thought of as $t(\gamma(\lambda))=\lambda$; precisely $a=1$, $b=0$) \ts{10}{1}{49:21}. So all the $a$-factors are fixed once and for all.
As it stands the equation is correct, but one has to think of it as evaluated along the curve: the vector fields $W_i$ are evaluated where the particle passes; one usually does not write this, to keep the equation tidy \ts{10}{1}{50:47}.

\subsection{The second law in a chart}\label{ssec:nd-chart}

That is it already; now quickly in a chart: nothing really new happens, but it is good to be sure that the curve must still be at least affinely parametrised, so that we may make this choice at all \ts{10}{1}{51:39}. The situation is in a chart $(U,x)$ of a stratified atlas (as always). We saw that the interaction fields are required to be spatial, so
writing the equation yields in the zeroth component a $q^0$: let $q=x\circ\gamma$ be the chart image of the curve \ts{10}{1}{52:20}. There is (no assumption yet that it has to do with $t$)
\[
  m\Bigl(\ddot q^{\,0}(\lambda)+\Gamma^0{}_{mn}\bigl(\gamma(\lambda)\bigr)\,\dot q^m(\lambda)\,\dot q^n(\lambda)\Bigr)\;=\;\sum_ic_i\,W_i^0\bigl(\gamma(\lambda)\bigr).
\]
The right side is the zeroth component of the resulting vector at the point where the particle passes, and \emph{zero}, since every $W_i$ is spatial; the $\Gamma^0{}_{mn}$ vanish because of the stratified atlas \ts{10}{1}{53:04}. So even after the introduction of the interaction fields
\begin{equation}\label{eq:nd-q0}
  \ddot q^{\,0}(\lambda)=0\quad\Longrightarrow\quad q^0(\lambda)=t\bigl(\gamma(\lambda)\bigr)=a\lambda+b ,
\end{equation}
as before; and this shows that we can indeed choose $a=1$ and $b=0$: the choice we made at the back is justified, because the interaction field was required to be spatial \ts{10}{1}{54:07}. (The lecturer notes that on the board the mass $m$ must stand in front of everything on the left; he had forgotten it \ts{10}{1}{55:10}.)

For the remaining spatial components $\alpha=1,2,3$ the double dot can be written, with the choice made, as the derivative with respect to $t$, and the terms with $\dot q^0=1$ give
\begin{equation}\label{eq:nd-chart}
  m\Bigl(\ddot q^{\,\alpha}(t)+\Gamma^\alpha{}_{00}\bigl(\gamma(t)\bigr)+2\,\Gamma^\alpha{}_{0\beta}\bigl(\gamma(t)\bigr)\,\dot q^{\,\beta}(t)+\Gamma^\alpha{}_{\beta\gamma}\bigl(\gamma(t)\bigr)\,\dot q^{\,\beta}(t)\,\dot q^{\,\gamma}(t)\Bigr)\;=\;\sum_{i=1}^Lc_i\,W_i^{\alpha}\bigl(\gamma(t)\bigr)
\end{equation}
(on the board the lecture had first written the mixed term without the factor $2$ and without the last term, and a student pointed out the omission of the last term, see below; the form \eqref{eq:nd-chart} is the complete one, with the factor $2$ from the symmetry $\Gamma^\alpha{}_{0\beta}=\Gamma^\alpha{}_{\beta0}$) \ts{10}{1}{55:30}--\ts{10}{1}{56:13}.

\paragraph{Special case: inertial system.} In an inertial system we computed from the abstract definition (\cref{def:nk-inertial}) that in a stratified atlas $\Gamma^\alpha{}_{00}$ and $\Gamma^\alpha{}_{0\beta}$ vanish \ts{10}{1}{56:33}. What remains is
$m\,\ddot q^{\,\alpha}(t)=F^\alpha(\gamma(t))$, where for the sum on the right the lecturer now simply introduces the abbreviation $F^\alpha$ for the force at the point where the particle passes: ``you see, this is simply the completely naively written Newton equation, in a spacetime in which you have a global
inertial system at all, and then also in coordinates in which you are actually in an inertial system'' (there can be an inertial system and you sit in the carousel and are no longer in it; one has to go into this inertial system) \ts{10}{1}{57:14}.

\begin{coursenote}[A student: the term $\Gamma^\alpha{}_{\beta\gamma}\dot q^\beta\dot q^\gamma$ is missing (lecture error, corrected in class)]\label{cn:nd-missing}
In the chart computation the lecturer had left out the term $\Gamma^\alpha{}_{\beta\gamma}\dot q^\beta\dot q^\gamma$ on the left-hand side (``a bracket'', ``my special bracket'': a student asked for it, ``thank you, I had omitted it, excellent''). It is the last term of \eqref{eq:nd-chart}.
It does \emph{not} vanish in an inertial system (the non-acceleration and non-rotation conditions of \cref{def:nk-inertial} concern only the $\Gamma^\alpha{}_{00}$ and $\Gamma^\alpha{}_{0\beta}$), so to get the school form $m\ddot q^\alpha=F^\alpha$ one needs a further assumption: that the whole Newtonian spacetime
is \emph{flat} (\cref{def:nk-flat}). The special case is therefore: \textbf{inertial system and flat Newtonian spacetime} \ts{10}{1}{59:42}--\ts{10}{1}{1:00:24}. A student then remarks on the index $\gamma$ used both for the curve and as a summation index: the lecturer agrees that it is an unlucky choice
(``only a madman would confuse it'') but keeps it, with the remark that in the formula $\gamma$ stands as the curve in $\gamma(t)$ and as an index on the $\Gamma$'s \ts{10}{1}{1:00:49}.
\end{coursenote}
\begin{theorem}[Second law in an inertial system of a flat Newtonian spacetime]\label{thm:nd-school}
In a chart of an inertial system of a flat Newtonian spacetime (\cref{def:nk-flat}), the world line $q(t)=(q^1(t),q^2(t),q^3(t))$ of a particle satisfies
\begin{equation}\label{eq:nd-school}
  m\,\ddot q^{\,\alpha}(t)\;=\;F^\alpha\bigl(\gamma(t)\bigr),\qquad\alpha=1,2,3 .
\end{equation}
\end{theorem}
This is the equation you know; if I had written it on the first day, you would have said ``yes, that is the whole story''; but it is by no means the whole story \ts{10}{1}{58:17}. If one starts from \eqref{eq:nd-school} and wants to go backwards, to the general case, \emph{one can never derive the general case from the special case}:
the way from there to the general one, with the inspiration of Maxwell theory, is Einstein's great achievement \ts{10}{1}{58:37}. Why should we start today with something which we know to be a huge special case? We would only confuse you; therefore we began as we did. It does give confidence: in an inertial system,
if there is one, one may do what we computed in experimental physics; it is completely right, hence sensible things come out, but with a whole lot of background which you hopefully now understand \ts{10}{1}{58:58}.

\paragraph{On the problem sheet.} (A small correction to the problem sheet: the hand-in is of course only on Tuesday, 3 June, two days earlier \ts{10}{1}{1:01:10}.) ``On this sheet there are a few pearls of physics'' \ts{10}{1}{1:01:30}:
\begin{itemize}
  \item For the first exercise there is a whole ``sermon'' at the bottom. For all who are worried that \emph{Lagrangian mechanics} comes up short: you have done it all already, you just did not notice; the lecturer will explain it in several extra lectures devoted to it, but we have already done almost everything we need there \ts{10}{1}{1:01:50}.
        The previous sheet already contained the variational-principles exercises: ``damn important'', go to the exercise group today and look at them: how it works with these \emph{variations} is very important; he will explain it again in lecture, but one never understands it as well as when one has tried it before \ts{10}{1}{1:02:32}.
  \item Then the flat Newtonian spacetime: it begins abstractly, ``show that there is still such a flat stratified atlas'', a very nice exercise in linear algebra \ts{10}{1}{1:02:53}.
  \item Wonderful is exercise 2e: change to a \emph{rotating coordinate system}, with the technique of the dreaded exponential function, $e^{\text{matrix}}$ (``the meditating Buddha tells you how $e^{\text{matrix}}$ goes, it is there''). One computes the $\Gamma$'s: remember that in the old system all $\Gamma$'s are zero,
        so only the inhomogeneous term occurs; the end result is ``a whole lot of clutter''; in it one is to identify the term that is the \textbf{centrifugal acceleration} and the term that is the \textbf{Coriolis acceleration} \ts{10}{1}{1:03:14}. One knows the formulas, but they are written not in the tensor calculus but in the
        profane three-dimensional flat-space vector/pseudovector calculus with the \emph{cross product}; the connection is the totally antisymmetric $\varepsilon$ ($\varepsilon_{123}=1$, antisymmetric, zero otherwise) and the rule for the cross product of three-dimensional vectors from school, used backwards: in the end one has the centrifugal and the Coriolis acceleration as they appear in a rotating system with angular velocity vector $\vec\omega$
        (the lecturer's remark ``the very relaxed young lady down here'' refers to the figure on the sheet) \ts{10}{1}{1:03:57}.
  \item And the third, ``finally, finally, finally'': one recognises the right-hand side; if the constant is chosen correctly it is Newton's gravitational law of a point mass. One does \emph{not} let this gravitational force stay a force: one writes it as one of the $\Gamma$'s that would have to be zero if there were an inertial system,
        namely $\Gamma^\alpha{}_{00}$, and this is the exercise of the \textbf{curved Newtonian spacetime}: at the end one is also to find a few geodesics (the lecturer corrects himself: autoparallels) in this spacetime. These would be spirals; projected to the absolute space they are circles; and if one could solve the elliptic integrals (``of course no one can'') one could
        show that the ellipses work too: in the spacetime the spiral over such a circle is the planetary motion resolved in time, and in the spacetime this spiral \emph{is} an autoparallel, ``exactly as we wanted''. In the next lecture gravitation is put under the magnifying glass separately and he shows how to get all orbits, i.e.\ how the Kepler problem is solved \ts{10}{1}{1:04:57}.
\end{itemize}
(``That was the commercial for the problem sheet'' \ts{10}{1}{1:06:21}.)

% ==================================================================
\section{Work and potential energy}\label{sec:nd-work}

We go on with two extremely important notions: \textbf{work} (Arbeit) and \textbf{potential energy} (potenzielle Energie) \ts{10}{1}{1:06:44}. The first fact: work and potential energy \emph{cannot be defined on a mere Newtonian spacetime}; we need a spatial metric. What we need is a
\emph{Newtonian spacetime with absolute space}, the whole structure $(M,\mathcal O,\mathcal A,\nabla,t,G)$ as defined in \cref{def:as-nrz-as} \ts{10}{1}{1:07:05}. And we must distinguish cleanly: if one has this and interaction fields and looks at particles, the \emph{work} done on a particle when one moves it along a certain
path (made precise in a moment, in spacetime if you like) is always defined, but only under special circumstances is the \emph{potential energy} defined \ts{10}{1}{1:08:09}.

\begin{definition}[New concepts: work done on a particle (Arbeit)]\label{def:nd-work}
Let $(M,\mathcal O,\mathcal A,\nabla,t,G)$ be a Newtonian spacetime with absolute space, and work in a chart of an inertial system throughout this section (``this entire subchapter 3''; otherwise it makes no sense). Let $\sigma\colon[a,b]\to M$ be a
\textbf{future-pointing} (resp.\ \textbf{past-pointing}) smooth curve in the spacetime, on a closed parameter interval, \emph{not} an arbitrary parameter range \ts{10}{1}{1:08:35}. For a particle $(\gamma,m,c_1,\dots,c_L)$, in the presence of interaction fields $W_1,\dots,W_L$ (with force $F=\sum c_iW_i$), the number
\begin{equation}\label{eq:nd-work}
  \mathcal W_{\sigma}\;:=\;-\,\epsilon\int_a^b\mathrm d\lambda\;G_{ab}\bigl(\sigma(\lambda)\bigr)\,F^{a}\bigl(\sigma(\lambda)\bigr)\,\dot\sigma^{b}(\lambda)
  \;=\;-\,\epsilon\int_a^b\mathrm d\lambda\;g_{\alpha\beta}\bigl(\sigma(\lambda)\bigr)\,F^{\alpha}\bigl(\sigma(\lambda)\bigr)\,\dot\sigma^{\beta}(\lambda),
\end{equation}
where $\dot\sigma^b$ are the components of the tangent vector of $\sigma$ in the chart, $\epsilon=+1$ for a future-pointing and $\epsilon=-1$ for a past-pointing curve, is called the \textbf{work done on the particle along the spacetime curve $\sigma$} (arbeit am Teilchen verrichtet) \ts{10}{1}{1:09:48}--\ts{10}{1}{1:14:02}.
\end{definition}
(The board has the notation $\mathcal W_{\sigma}$, ``the number depends on the choice of your curve $\sigma$; $\sigma(a)\to\sigma(b)$, a bit of overkill''; the placement of $\epsilon$: it is the sign, stated at \ts{10}{1}{1:14:44}, ``this epsilon stands here below''; whether in front of the integral or in the exponent of an orientation convention is not recoverable from the captions, [?]; a student remarked at \ts{10}{1}{1:37:44} on a missing sign, and at \ts{10}{1}{1:27:59} on the omitted $\epsilon$, which the lecturer thanks for.)
Observe that $\sigma$ is \emph{not necessarily} the world line of the particle, and in general not at all identical with it: any curve. We use the spatial spacetime metric (we have it, since we took a Newtonian spacetime with absolute space), evaluated along $\sigma$; it eats the force field $F^a$ (with an upper index, evaluated again where $\sigma$ passes;
the force field is the weighted sum with the charges relevant for this particle of the interaction fields) and the tangent vector of $\sigma$, $\beta$-component. Types: $G$ is a spacetime $(0,2)$ tensor field, $F$ a $(1,0)$ tensor field, the tangent a $(1,0)$ field (at least where I pass); this is a contraction, and contraction reduces the rank, so what stands there is a spacetime $(0,0)$ tensor field, a function on the spacetime, but evaluated only along $\sigma$; one can then integrate over the curve parameter and gets a number \ts{10}{1}{1:11:54}.
That the work is done \emph{on} the particle and not \emph{by} the particle is what the minus sign expresses \ts{10}{1}{1:14:02}.

\begin{intuition}[Work need not really happen: moving the particle in thought]\label{int:nd-work}
How can $\sigma$ differ from the curve of the particle, and what sense does it make? The point: the work need not really be done for me to think about it \ts{10}{1}{1:15:06}. Suppose a particle (free, or already on its world line because a force acts, no matter) satisfies the second axiom. Of the interaction fields only their values where the particle passes are relevant for its motion.
But one can consider another curve $\sigma$, for instance starting ``in the middle of the thicket'' at $\sigma(a)$ and ending at $\sigma(b)$ (a thought experiment): in thought one can \emph{force} the particle there, while it would, if let go in the interaction fields, simply follow its own world line \ts{10}{1}{1:16:08}. Example: an electron in an electric field. If I let it go, one knows what it does.
But one can take exactly this electron, grab it and pull it, with effort, backwards against the electric field; or let it run and merely imagine that one pulls it: this is the curve $\sigma$ \ts{10}{1}{1:17:51}. So one does not necessarily realise this work, but one can talk about it. If one wants to go backwards along the curve, the tangent vector is then no longer future-pointing but past-pointing, and the
sign change is caught by the $\epsilon$ (``very academic, just so that it fits in all cases and one gets no wrong ideas''). Evaluating the integral gives exactly this quantity, the work (maybe only in thought) done on the particle when I pull it along this path \ts{10}{1}{1:18:32}.
\end{intuition}
\begin{coursenote}[Questions about what enters $\mathcal W_\sigma$]
A student asks whether the interaction (the coupling of the particle to the field) matters: yes, ``that is right, the interaction is in it; it makes a difference whether a particle couples to something or not: in this $F$ are the $c$'s, how they couple to the $W$''; $\gamma$ and $m$ play no role (``right''). Then: ``but what about the gravitational field? There $m$ must play a role'';
answer: gravitation is not a force for us \ts{10}{1}{1:19:15}. Another question about the curve $\sigma$ being continuous: for the definition one needs the tangent vector, so $C^1$ would be enough; in general ``curve'' means smooth curve; one could sharpen all theories and definitions (here a $C^2$ tensor field, there only $C^1$, or a $C^3$ manifold instead of $C^\infty$), but whenever one has a $C^k$ manifold one can make
a $C^\infty$ manifold out of it; one says explicitly what is needed only for the structure of the spacetime, which is conceptually important at the moment \ts{10}{1}{1:31:27}.
\end{coursenote}

\subsection{Evaluation of the work}\label{ssec:nd-work-eval}

Now an immediate observation. Since $G$ is a \emph{spatial} spacetime metric and also the vector field $F$ is spatial: you remember that if one has a spatial spacetime vector field and writes two arbitrary spacetime vector fields into $G$, only the spatial components contribute (the lecture of \cref{ch:c09}, \eqref{eq:as-G-projection}).
So \eqref{eq:nd-work} reduces to the components $g_{\alpha\beta}$ with $\delta^\alpha{}_a\delta^\beta{}_b$, and what remains is $g_{\alpha\beta}F^\alpha\dot\sigma^\beta$ with the spatial component of the tangent vector (in a chart, the $\beta$-th component of the derivative of the chart image of $\sigma$ with respect to $\lambda$) \ts{10}{1}{1:20:02}.
There is no need to write the \emph{time} here as parameter: one could parametrise $\sigma$ by time but need not; one checks easily that under a change of parameter the tangent vector changes by the inverse factor, so the number $\mathcal W$ is completely independent of how fast, or with which parameter, one runs along $\sigma$: it is \textbf{reparametrisation invariant} \ts{10}{1}{1:21:46}.
Hence the work is computed in the form familiar to you, projecting onto the space:
\begin{equation}\label{eq:nd-work-chart}
  \mathcal W_\sigma\;=\;-\,\epsilon\int_a^b\mathrm d\lambda\;g_{\alpha\beta}\bigl(\sigma(\lambda)\bigr)\,F^\alpha\bigl(\sigma(\lambda)\bigr)\,\dot\sigma^\beta(\lambda) ,
\end{equation}
the curve in the chart differentiated. Often one writes $\mathrm d\lambda\,\dot\sigma$ as $\mathrm d\sigma$, as when one writes in mechanics $\mathrm d\vec r$ instead of $\vec v\,\mathrm dt$, and the work integral becomes $\int\mathrm d\vec r\cdot\vec F$ ``with all kinds of assumptions''; ``that is precisely the work formula that you know'' \ts{10}{1}{1:22:07}.
(A student remarks that $G$ in the line before should be the spatial spacetime metric: right, ``that is a spatial spacetime metric'', thank you \ts{10}{1}{1:23:30}.)

\paragraph{The lesson: what is hidden in the work formula.} The work in this form can be defined only on a Newtonian spacetime with absolute space, i.e.\ with inertial system and so on. \emph{Not} necessarily flat: the $g_{\alpha\beta}$ could be non-trivial, giving non-trivial $\Gamma^\alpha{}_{\beta\gamma}$, the \emph{spaces could still be curved} \ts{10}{1}{1:24:13}.

\begin{example}[The rocket and the three-dimensional sphere (Columbus problem)]\label{ex:nd-rocket}
It may seem odd, but suppose you sit in a rocket, the world is Newtonian, and you fly off in one direction. It could be that the universe is such that, if you fly long enough, you arrive back here, but \emph{from the other side}: ``the Columbus problem''. One sails and sails and ``hup, I am back''. Just as one thought the Earth a disc and was then surprised (``this was known in antiquity, people were better''), one does not naively
assume that going on in one direction one comes back from behind; but we know that it is so on the Earth (without ocean and such). The universe could be like that for three-dimensional space too: one flies off, and after 18 billion light years or so comes back from the other side. If so, space is a sphere, a three-dimensional sphere (``I cannot draw one, a dimension higher''), and then we need the $g$'s to be \emph{non-trivial}
(and this has nothing to do with the time). \ts{10}{1}{1:24:34}
\end{example}
So it is not an academic exercise that $g$ stands there and is not fixed to the flat one (one writes inner products as ``$1,1,1$'' and such, ``you can do that, but you rob yourself of every insight that one can have''). We theorists are responsible for what is possible and conceivable; ``so far it has often turned out in physics that what was conceivable was far from innovative enough compared with what actually was''.
This prepares for everything; we could do, say, Newtonian cosmology with it \ts{10}{1}{1:25:56}. A student remarks that the formula no longer depends on the change in time, so one can imagine the situation projected into one of these spaces and go only from the point here to the point there: ``that is right, exactly; that is the usual picture, we have the space and in space we take a path integral over the force
projected on the tangent vector of the path, and then we get the work'' \ts{10}{1}{1:26:58}. But look at which conditions are hidden in it: we are so general that our universe, even with the rocket, still fits; ``you can never get that in if you define the work with the vector arrow and so on: then you are always in flat space, and the global structure of the universe may be quite different'' \ts{10}{1}{1:27:18}.
(An omitted $\epsilon$ in the formula is noted by a student and corrected \ts{10}{1}{1:27:59}.)

\paragraph{The work depends on the path.} One can always compute the work, but it depends on the path $\sigma$, obviously from the definition; and in general, for arbitrary vector interaction fields, it is \emph{not} the same whether I choose one path from the starting to the end point or another one, both everywhere future-pointing \ts{10}{1}{1:28:19}.
\begin{proposition}[The work in general depends on the path]\label{prop:nd-pathdep}
Let $\sigma\colon[a,b]\to M$ and $\tilde\sigma\colon[c,d]\to M$ be two smooth curves as in \cref{def:nd-work} with $\sigma(a)=\tilde\sigma(c)$ and $\sigma(b)=\tilde\sigma(d)$ (the same initial and end point in the spacetime). Then in general, for general unrestricted interaction fields $W_1,\dots,W_L$,
\[
  \mathcal W_\sigma\;\neq\;\mathcal W_{\tilde\sigma}
\]
for the work done on the same particle along $\sigma$ and along $\tilde\sigma$ (for the work the world line of the particle plays no role) \ts{10}{1}{1:29:39}--\ts{10}{1}{1:31:06}.
\end{proposition}
(``In general'' here means for general, not further restricted interaction fields; the lecture states it as a warning, not as a theorem with a proof; the examples are the usual non-conservative fields.)

\subsection{Conservative interaction fields and the potential energy}\label{ssec:nd-potential}

There is, however, a special class of interaction fields for which this is nevertheless the case \ts{10}{1}{1:32:50}.
\begin{definition}[New concepts: conservative interaction field (konservatives Wechselwirkungsfeld)]\label{def:nd-conservative}
An interaction field $W$ (a spatial spacetime vector field, considered individually; say the $i$-th one, $W_i$) is called \textbf{conservative} if the integral
\begin{equation}\label{eq:nd-conservative}
  \mathcal W^{(i)}_\sigma\;:=\;-\,\epsilon\int_a^b\mathrm d\lambda\;g_{\alpha\beta}\bigl(\sigma(\lambda)\bigr)\,W_i^\alpha\bigl(\sigma(\lambda)\bigr)\,\dot\sigma^\beta(\lambda)
\end{equation}
(the part of the work coming from the $i$-th interaction field only, without the charge $c_i$) is \emph{independent of $\sigma$}: that is, for fixed endpoints $\sigma(a)$ and $\sigma(b)$ (one must not change them too) and all other future-pointing curves between the same endpoints, as in the previous definition, it takes the same value \ts{10}{1}{1:33:19}--\ts{10}{1}{1:35:33}.
\end{definition}
It follows: if a particle $(\gamma,m,c_1,\dots,c_L)$ couples \emph{only to conservative} interaction fields, then the work done on the particle is also only dependent on the initial and the end point of the curve $\sigma$ \ts{10}{1}{1:36:15}. (There may be other, non-conservative, interaction fields; but if the charges coupling to them are zero for a concrete particle, the particle feels nothing of them: its work is $\sum_ic_i\mathcal W^{(i)}_\sigma$,
a sum over the conservative ones; the entire force acting on the particle has the same property \ts{10}{1}{1:37:03}.) A student notices a missing sign ``$-Y$'' on the board, which the lecturer corrects \ts{10}{1}{1:37:44}.

Picture: one draws the plane (not really needed, but one always draws it, ``so that you know that we are in spacetime''). If one looks at the particle and it is irrelevant whether it goes here or there, with $\sigma$ and the other $\tilde\sigma$, the situation is always the same \ts{10}{1}{1:38:07}. Then one defines the potential energy:
\begin{definition}[New concepts: potential energy (potenzielle Energie)]\label{def:nd-potential}
In exactly this case (the particle couples only to conservative interaction fields), the \textbf{potential energy of the particle with respect to the point $Q$} is the function on the spacetime
\begin{equation}\label{eq:nd-potential}
  V(P)\;:=\;\mathcal W_\sigma\qquad\text{for any future-pointing curve }\sigma\text{ with }\sigma(a)=Q,\ \sigma(b)=P ,
\end{equation}
where $P,Q$ are points in the spacetime \ts{10}{1}{1:38:28}--\ts{10}{1}{1:39:36}.
\end{definition}
(The value does not depend on $\sigma$, which is why only $Q$ and $P$ appear on the right-hand side: $P$ is what remains as the argument.) You know that when one talks about the potential energy of this piece of chalk, $mgh$, one has to say where $h=0$ is: two metres lower, up there: one chooses the reference point $Q$. And then one can think of a particle, and of what potential
energy it would have over there \emph{without bringing it there}; if one wants to determine the potential energy experimentally one must really bring it there (then $\sigma$ would coincide with $\gamma$), but one can also just think it; ``that is my point, and that is the potential energy'' \ts{10}{1}{1:39:59}. Once the reference point is fixed, it depends only on where the particle is.
``But how do I get there? Never mind'': it can be never-mind only if the interaction fields to which the particle couples are conservative; otherwise \textbf{there is no potential energy} \ts{10}{1}{1:40:40}.

\begin{coursenote}[There is no \emph{the} potential energy]\label{cn:nd-nopot}
The statements ``there is a kinetic energy and there is a potential energy'' are ``complete nonsense'' as general statements: a potential energy exists only in a Newtonian spacetime with absolute space and with interaction fields that are conservative, at least those to which the concrete particle couples. If, say, half of the interaction fields are conservative and
the other half not, then there are particles for which there is a potential energy and others for which there is none, namely those with non-vanishing charges coupling to non-conservative fields \ts{10}{1}{1:41:01}. The gravitational field is conservative in Newtonian mechanics (by the field equation, which we look at next time), and \emph{that} is why there is a gravitational potential.
There is the distinction \emph{potential} versus \emph{potential energy} (the latter is the former times the charge, or not); but this exists really only for conservative force fields \ts{10}{1}{1:42:05}. ``Again our conceptual pedantry shows differences that one could otherwise only vaguely surmise or discuss case by case: now all assumptions lie explicitly on the table, and there should be no confusion'' \ts{10}{1}{1:42:46}.
\end{coursenote}
\begin{exercise}[Recovering the force from the potential energy (announced for the next problem sheet)]\label{ex:nd-grad}
On the next problem sheet one is to show that a conservative force field can be \emph{reconstructed} from the potential energy: ``minus the gradient''. The gradient is a covector (in a chart, the partial derivative, index down), so to get a force \emph{vector} one needs the spatial spacetime metric again, to raise the index; more structure than is usually written when everything is assumed flat; also the sufficient condition that
the rotation of the vector field vanishes needs the metric to compute the rotation. Conceptually one may not ignore this; in Newtonian cosmology it would make a big difference that the $g$ stands there \ts{10}{1}{1:43:07}.
\end{exercise}
Next time: how the gravitational force really works, and we compute all planetary orbits \ts{10}{1}{1:44:31}.

% ==================================================================
\section*{Key concepts}
\begin{keyconcepts}
\textbf{Retrieval questions.}
\begin{enumerate}
  \item Why is the force concept in Newton strange (instantaneous action), and what is Faraday's field idea?
  \item Define a spatial spacetime vector field and an interaction field; how many interaction fields are there, and how does the number depend on the universe?
  \item Define the force on a particle. Why does it matter that the interaction fields are independent of the particle, while the charges are not?
  \item Why does the first axiom speak of vanishing charges instead of vanishing force?
  \item Which fundamental interactions are classically describable, and why does electromagnetism not fit the Newtonian spacetime? What is an effective interaction (two examples)?
  \item State the second Newtonian axiom precisely; why can one choose $\lambda=t$? Write it in a chart of a stratified atlas and say which terms vanish in an inertial system and which extra assumption gives $m\ddot q^\alpha=F^\alpha$.
  \item Define the work done on a particle along a spacetime curve $\sigma$. Why does it need the absolute space, and why does it not depend on the parametrisation or on the world line of the particle?
  \item Define a conservative interaction field and the potential energy of a particle with respect to $Q$. When does it not exist?
\end{enumerate}
\textbf{Concepts.}
\begin{itemize}
  \item Interaction field and force (\cref{def:nd-field,def:nd-force}, \eqref{eq:nd-force}); fundamental versus effective interactions (\cref{ssec:nd-fundamental,ssec:nd-effective}).
  \item Second law (\cref{ax:nd-second}, \eqref{eq:nd-second}); chart form \eqref{eq:nd-chart}; school form (\cref{thm:nd-school}, \cref{cn:nd-missing}).
  \item Work (\cref{def:nd-work}, \eqref{eq:nd-work-chart}); path dependence (\cref{prop:nd-pathdep}).
  \item Conservative field and potential energy (\cref{def:nd-conservative,def:nd-potential}); \cref{cn:nd-nopot}.
\end{itemize}
\end{keyconcepts}

\flushside
\end{paracol}
